\documentclass[11pt]{article}
\usepackage[margin=1in]{geometry}
\usepackage{amsmath}
\usepackage{booktabs}
\usepackage{hyperref}

\title{CS 610 Coding Assignment 4 --- Report}
\author{}
\date{}

\begin{document}
\maketitle

\section*{Environment}
\begin{itemize}
  \item \textbf{GPU:} NVIDIA A100 (Ampere, CC 8.0, 108 SMs, HBM2e)
  \item \textbf{Compile flags:} \texttt{nvcc -O2 -arch=sm\_80}
\end{itemize}

\section*{Performance Results}

\begin{table}[h]
\centering
\caption{Execution times for all four reduction methods (A100).}
\begin{tabular}{lcc}
\toprule
\textbf{Method} & \textbf{Student.dat (2048)} & \textbf{Student\_large.dat ($2^{20}$)} \\
\midrule
\texttt{maximumMark\_atomic}     & 4.95 ms        & $\sim$208{,}953 ms \\
\texttt{maximumMark\_recursive}  & 0.063 ms       & 6.58 ms            \\
\texttt{maximumMark\_SM}         & 0.034 ms       & 2.08 ms            \\
\texttt{maximumMark\_shuffle}    & 0.036 ms       & 2.11 ms            \\
\bottomrule
\end{tabular}
\end{table}

\noindent All four methods passed the \texttt{cpuValidation()} function, which computes the maximum
sequentially on the CPU and compares against each GPU result:
\begin{itemize}
  \item \textbf{Student.dat}: All four methods matched CPU exactly (mark = 99.993896, student = 32651).
  \item \textbf{Student\_large.dat}: Recursive, SM, and Shuffle matched CPU exactly (mark = 100.0, student = 25663).
        Atomics found the correct mark (100.0) but a different student ID (22760 vs.\ 25663)
        due to non-deterministic tie-breaking among multiple students with the same maximum score.
\end{itemize}

\section*{Cross-GPU Comparison}

For reference, Table~2 shows the same kernels executed on an RTX 4090 (Ada Lovelace, CC 8.9, 128 SMs, GDDR6X).

\begin{table}[h]
\centering
\caption{Execution times on RTX 4090 for comparison.}
\begin{tabular}{lcc}
\toprule
\textbf{Method} & \textbf{Student.dat (2048)} & \textbf{Student\_large.dat ($2^{20}$)} \\
\midrule
\texttt{maximumMark\_atomic}     & 2.55 ms      & $\sim$81{,}031 ms \\
\texttt{maximumMark\_recursive}  & 0.054 ms     & 4.93 ms           \\
\texttt{maximumMark\_SM}         & 0.025 ms     & 1.33 ms           \\
\texttt{maximumMark\_shuffle}    & 0.026 ms     & 1.31 ms           \\
\bottomrule
\end{tabular}
\end{table}

\noindent The RTX 4090 was faster across all methods despite the A100 having $2\times$ higher memory bandwidth (HBM2e at $\sim$2~TB/s vs.\ GDDR6X at $\sim$1~TB/s).
This is because parallel reduction on these dataset sizes is not purely memory-bandwidth-bound---the RTX 4090's
higher SM count (128 vs.\ 108) and $2\times$ higher FP32 throughput per SM (Ada's dual-issue datapath) give it
an advantage on compute-light kernels where the working set fits in the large L2 cache ($\sim$72~MB vs.\ 40~MB).
The atomic approach shows the largest gap ($2\times$ slower on A100) because the A100 supports more concurrent warps
per SM (64 vs.\ 48), which increases contention on the spin-lock.

\section*{Analysis}

\textbf{Atomic approach.}
The spin-lock serializes all threads through a single critical section, making it effectively $O(N)$ serial work.
On the A100, scaling from 2{,}048 to $2^{20}$ records increased execution time by $\sim$42{,}200$\times$, far worse than the 512$\times$ increase in data size, because contention on the lock variable grows super-linearly as more threads compete for the same cache line.

\textbf{Recursive approach.}
Each kernel launch halves the data via pairwise comparison in shared memory.
This requires $\lceil\log_2(N / \text{TPB})\rceil$ sequential kernel launches (3 for 2{,}048 records, 12 for $2^{20}$).
Kernel launch overhead dominates at small $N$; at large $N$ it is fast but not optimal because only half the threads do useful work per launch and each launch incurs synchronization overhead.

\textbf{Shared memory approach.}
A single kernel launch fully reduces each block of 256 threads using tree-based parallel reduction ($\log_2(256) = 8$ steps, all in shared memory).
This is the most efficient because it minimizes kernel launches (just one) and keeps all reduction traffic in fast shared memory.
The CPU only needs to reduce $N / 256$ block-level results.

\textbf{Shuffle approach.}
Warp-level reduction using \texttt{\_\_shfl\_down\_sync} avoids shared memory entirely; values are exchanged directly through registers via the warp shuffle unit.
Performance is nearly identical to the shared memory approach.
The shuffle kernel leaves more partial results for the CPU ($N/32$ warp results vs.\ $N/256$ block results), but this overhead is negligible.

\end{document}
